problem:  
Let \( G = \mathrm{SU}(3) \) act on an 8‑dimensional, compact, simply connected manifold \( M^8 \) with cohomogeneity one.  
Suppose that:  
- the **principal orbits** are Aloff–Wallach spaces \( W_{p,q} = \mathrm{SU}(3)/\mathrm{U}(1)_{p,q} \), where \( (p,q) \) are coprime integers,  
- the **singular orbits** correspond to closed subgroups \( K_-, K_+ \subset G \) containing \( \mathrm{U}(1)_{p,q} \), and  
- \( M \) admits an \( \mathrm{SU}(3) \)-invariant Riemannian metric of **positive sectional curvature**.  

Open question:  
For fixed singular isotropy types \( K_-, K_+ \), does the diffeomorphism type of \( M \) depend only on the pair \((p,q)\)?  
Equivalently, are there at most finitely many diffeomorphism types of positively curved cohomogeneity‑one 8‑manifolds with orbit diagram  
\[
(\mathrm{SU}(3), K_-, K_+, \mathrm{U}(1)_{p,q})
\]
as \((p,q)\) varies among coprime integer pairs?

Why is it a "good" problem:  
This formulation is precise, dimensionally consistent, and directly positioned at the frontier of the classification of positively curved manifolds. It asks whether curvature and symmetry together impose topological rigidity on cohomogeneity‑one spaces built from Aloff–Wallach orbits.  

This problem lies at a crossroads among several deep areas of geometry and topology:  

1. **Positive curvature classification** – Extending beyond known examples (e.g., Aloff–Wallach and Eschenburg spaces), it asks whether new manifolds can arise or whether the structure locks to finitely many topological types.  

2. **Transformation group theory** – The group diagram \((G, K_-, K_+, H)\) encodes the full topological type of a cohomogeneity‑one manifold; investigating whether curvature can constrain this diagram ties curvature geometry to representation theory and Lie group topology.  

3. **Exotic structures and rigidity** – Aloff–Wallach spaces themselves include distinct homotopy types and potentially exotic smooth structures. Whether such differences affect (or are suppressed by) the ambient curvature constraints probes the delicate interface between Riemannian geometry and differential topology.  

The statement is concise yet leads into deep, unsolved territory—it invites geometric, topological, and analytic techniques for progress. Its combination of clarity, open status, and cross‑disciplinary depth makes it a genuinely **good** research-level problem.